\ifdefined\gkver\else\def\gkver{student}\fi
\documentclass[\gkver]{gaokaozhenti}
\begin{document}

\chapter{2008年四川理综}
%% paperType: 理综物理部分

\begin{enumerate}
\item
%% number: 14
%% paperName: 2008年四川理综第 14 题 6 分
%% typeId: 1
%% type: 单选题
%% chapter: 气体性质
%% point: 热力学第一定律
%% method: 无
%% score: 6
%% degree: 850
%% duplicateId: 0
%% body:
下列说法正确的是\xzanswer{D}

\fourchoices[answer=D]
{物体吸收热量，其温度一定升高}
{热量只能从高温物体向低温物体传递}
{遵守热力学第一定律的过程一定能实现}
{做功和热传递是改变物体内能的两种方式}

%% answer: D
%% memo:
\memoanswer{
由热力学第一定律可知，做功与热传递可以改变物体的内能，D 正确；故物体吸收热量时，其内能不一定增大，A 错；由热力学第二定律可知，宏观的热现象有方向性，但若通过外界做功，热量也可以从低温物体传到高温物体，B、C 错。
}

\item
%% number: 15
%% paperName: 2008年四川理综第 15 题 6 分
%% typeId: 1
%% type: 单选题
%% chapter: 原子和原子核
%% point: 核聚变
%% method: 无
%% score: 6
%% degree: 850
%% duplicateId: 0
%% body:
下列说法正确的是\xzanswer{A}

\fourchoices[answer=A]
{$\gamma$ 射线在电场和磁场中都不会发生偏转}
{$\beta$ 射线比 $\alpha$ 射线更容易使气体电离}
{太阳辐射的能量主要来源于重核裂变}
{核反应堆产生的能量来自轻核聚变}

%% answer: A
%% memo:
\memoanswer{
$\gamma$ 射线中的 $\gamma$ 光子不带电，故在电场与磁场中都不会发生偏转，A 正确；$\alpha$ 粒子的特点是电离能力很强，B 错；太阳辐射的能量主要来源于轻核的聚变，C 错；核反应堆产生的能量是来自于重核的裂变，D 错。
}

\item
%% number: 16
%% paperName: 2008年四川理综第 16 题 6 分
%% typeId: 2
%% type: 多选题
%% chapter: 电路
%% point: 动态分析
%% method: 无
%% score: 6
%% degree: 650
%% duplicateId: 0
%% body:
如图，一理想变压器原线圈接入一交流电源，副线圈电路中 $R_{1}$、$R_{2}$、$R_{3}$ 和 $R_{4}$ 均为固定电阻，开关 S 是闭合的。$V_{1}$ 和 $V_{2}$ 为理想电压表，读数分别为 $U_{1}$ 和 $U_{2}$；$A_{1}$、$A_{2}$ 和 $A_{3}$ 为理想电流表，读数分别为 $I_{1}$、$I_{2}$ 和 $I_{3}$。现断开 S，$U_{1}$ 数值不变，下列推断中正确的是\xzanswer{BC}

\onepicture[w=6.0cm]{figs/16.png}

\fourchoices[answer=BC]
{$U_{2}$ 变小、$I_{3}$ 变小}
{$U_{2}$ 不变、$I_{3}$ 变大}
{$I_{1}$ 变小、$I_{2}$ 变小}
{$I_{1}$ 变大、$I_{2}$ 变大}

%% answer: BC
%% memo:
\memoanswer{
因为变压器的匝数与 $U_{1}$ 不变，所以 $U_{2}$ 与两电压表的示数均不变。当 S 断开时，因为负载电阻增大，故次级线圈中的电流 $I_{2}$ 减小，由于输入功率等于输出功率，所以 $I_{1}$ 也将减小，C 正确；因为 $R_{1}$ 的电压减小，故 $R_{2}$、$R_{3}$ 两端的电压将增大，$I_{3}$ 变大，B 正确。
}

\item
%% number: 17
%% paperName: 2008年四川理综第 17 题 6 分
%% typeId: 1
%% type: 单选题
%% chapter: 电磁感应
%% point: 楞次定律推广
%% method: 无
%% score: 6
%% degree: 650
%% duplicateId: 0
%% body:
在沿水平方向的匀强磁场中，有一圆形金属线圈可绕沿其直径的竖直轴自由转动。开始时线圈静止，线圈平面与磁场方向既不平行也不垂直，所成的锐角为 $\alpha$。在磁场开始增强后的一个极短时间内，线圈平面\xzanswer{B}

\fourchoices[answer=B]
{维持不动}
{将向使 $\alpha$ 减小的方向转动}
{将向使 $\alpha$ 增大的方向转动}
{将转动，因不知磁场方向，不能确定 $\alpha$ 会增大还是会减小}

%% answer: B
%% memo:
\memoanswer{
由楞次定律可知，当磁场开始增强时，线圈平面转动的效果是为了减小线圈磁通量的增加，而线圈平面与磁场间的夹角越小时，通过的磁通量越小，所以将向使 $\alpha$ 减小的方向转动。
}

\item
%% number: 18
%% paperName: 2008年四川理综第 18 题 6 分
%% typeId: 2
%% type: 多选题
%% chapter: 直线运动
%% point: 运动图象
%% method: 无
%% score: 6
%% degree: 650
%% duplicateId: 0
%% body:
一物体沿固定斜面从静止开始向下运动，经过时间 $t_{0}$ 滑至斜面底端。已知在物体运动过程中物体所受的摩擦力恒定。若用 $F$、$v$、$s$ 和 $E$ 分别表示该物体所受的合力、物体的速度、位移和机械能，则下列图象中可能正确的是\xzanswer{AD}

\fourchoices[answer=AD, ispicture=true, h=2.0cm]
{figs/18a.png}
{figs/18b.png}
{figs/18c.png}
{figs/18d.png}

%% answer: AD
%% memo:
\memoanswer{
物体在沿斜面向下滑动过程中，所受的合力为重力沿斜面向下的分力及摩擦力，故大小不变，A 正确；而物体在此合力作用下作匀加速运动，$v=at$，$s=\frac{1}{2}at^{2}$，所以 B、C 错；物体受摩擦力作用，总的机械能将减小，D 正确。
}

\item
%% number: 19
%% paperName: 2008年四川理综第 19 题 6 分
%% typeId: 2
%% type: 多选题
%% chapter: 机械波
%% point: 波与振动图象
%% method: 无
%% score: 6
%% degree: 450
%% duplicateId: 0
%% body:
一列简谐横波沿直线传播，该直线上的 a、b 两点相距 $4.42\Um$。图中实、虚两条曲线分别表示平衡位置在 a、b 两点处质点的振动曲线。从图示可知\xzanswer{AC}

\onepicture[w=7.5cm]{figs/19.png}

\fourchoices[answer=AC]
{此列波的频率一定是 $10\UHz$}
{此列波的波长一定是 $0.1\Um$}
{此列波的传播速度可能是 $34\Ums$}
{a 点一定比 b 点距波源近}

%% answer: AC
%% memo:
\memoanswer{
由振动图像可知，振动周期为 $0.1\Us$，故此列波的频率一定是 $10\UHz$，A 正确；因为波速为 $v=\frac{4.42}{0.1n-0.07}\Ums$，当 $n=2$ 时，波速为 $34\Ums$，故 C 正确；由图不能断定波长一定是 $0.1\Um$，也无法确定哪一点距波源近一些。
}

\item
%% number: 20
%% paperName: 2008年四川理综第 20 题 6 分
%% typeId: 1
%% type: 单选题
%% chapter: 曲线运动
%% point: 开普勒定律
%% method: 无
%% score: 6
%% degree: 650
%% duplicateId: 0
%% body:
1990 年 4 月 25 日，科学家将哈勃天文望远镜送上距地球表面约 $600\Ukm$ 的高空，使得人类对宇宙中星体的观测与研究有了极大的进展。假设哈勃望远镜沿圆轨道绕地球运行。已知地球半径为 $6.4\times10^{6}\Um$，利用地球同步卫星与地球表面的距离为 $3.6\times10^{7}\Um$ 这一事实可得到哈勃望远镜绕地球运行的周期。以下数据中最接近其运行周期的是\xzanswer{B}

\fourchoices[answer=B]
{0.6 小时}
{1.6 小时}
{4.0 小时}
{24 小时}

%% answer: B
%% memo:
\memoanswer{
由开普勒行星运动定律可知，$\frac{R^{3}}{T^{2}}=$ 恒量，所以 $\frac{(r+h_{1})^{3}}{t_{1}^{2}}=\frac{(r+h_{2})^{3}}{t_{2}^{2}}$，其中 $r$ 为地球的半径，$h_{1}$、$t_{1}$、$h_{2}$、$t_{2}$ 分别表示望远镜到地表的距离、望远镜的周期、同步卫星距地表的距离、同步卫星的周期（24 小时），代入数据得 $t_{1}=1.6$ 小时。
}

\item
%% number: 21
%% paperName: 2008年四川理综第 21 题 6 分
%% typeId: 1
%% type: 单选题
%% chapter: 光学
%% point: 光的折射
%% method: 无
%% score: 6
%% degree: 650
%% duplicateId: 0
%% body:
如图，一束单色光射入一玻璃球体，入射角为 $60^{\circ}$。已知光线在玻璃球内经一次反射后，再次折射回到空气中时与入射光线平行。此玻璃的折射率为\xzanswer{C}

\onepicture[w=4.0cm]{figs/21.png}

\fourchoices[answer=C]
{$\sqrt{2}$}
{1.5}
{$\sqrt{3}$}
{2}

%% answer: C
%% memo:
\memoanswer{
如图，A、C 为折射点，B 为反射点，作 OD 平行于入射光线，故 $\angle AOD=\angle COD=60^{\circ}$，所以 $\angle OAB=30^{\circ}$，玻璃的折射率 $n=\frac{\sin60^{\circ}}{\sin30^{\circ}}=\sqrt{3}$。

\onepicture[w=4.2cm]{figs/21_an.png}
}

\item
%% number: 22
%% paperName: 2008年四川理综第 22 题 8 分
%% typeId: 4
%% type: 实验题
%% chapter: 电路
%% point: 电路实验
%% method: 无
%% score: 8
%% degree: 650
%% duplicateId: 0
%% body:
\begin{enumerate}
	\item
	Ⅰ．（9 分）一水平放置的圆盘绕过其圆心的竖直轴匀速转动。盘边缘上固定一竖直的挡光片。盘转动时挡光片从一光电数字计时器的光电门的狭缝中经过，如图 1 所示。图 2 为光电数字计时器的示意图。光源 A 中射出的光可照到 B 中的接收器上。若 A、B 间的光路被遮断，显示器 C 上可显示出光线被遮住的时间。

	挡光片的宽度用螺旋测微器测得，结果如图 3 所示。圆盘直径用游标卡尺测得，结果如图 4 所示。由图可知，

	（1）挡光片的宽度为 \tkanswer[2cm]{10.243} $\Umm$。

	（2）圆盘的直径为 \tkanswer[2cm]{24.220} $\Ucm$。

	（3）若光电数字计时器所显示的时间为 $50.0\UmsT$，则圆盘转动的角速度为 \tkanswer[2cm]{1.69} $\Urads$（保留 3 位有效数字）。

	\onepicture[w=11cm]{figs/22a.png}

	\item
	Ⅱ．（8 分）图为用伏安法测量电阻的原理图。图中 V 为电压表，内阻为 $4000\UO$；mA 为电流表，内阻为 $50\UO$。$E$ 为电源，$R$ 为电阻箱，$R_{x}$ 为待测电阻，S 为开关。

	（1）当开关闭合后电压表读数 $U=1.6\UV$，电流表读数 $I=2.0\UmA$。若将 $R_{x}=\frac{U}{I}$ 作为测量值，所得结果的百分误差是 \tkanswer[2cm]{20\%}。

	（2）若将电流表改为内接。开关闭合后，重新测得电压表读数和电流表读数，仍将电压表读数与电流表读数之比作为测量值，这时结果的百分误差是 \tkanswer[2cm]{5\%}。（百分误差 $=\frac{\text{实际值}-\text{测量值}}{\text{实际值}}\times100\%$）

	\onepicture[w=4.0cm, align=right]{figs/22b.png}
\end{enumerate}

%% answer:
\jdanswer{
\begin{enumerate}
	\item
	Ⅰ．（1）$10.243\Umm$；（2）$24.220\Ucm$；（3）$1.69\Urads$。

	\item
	Ⅱ．（1）$20\%$；（2）$5\%$。
\end{enumerate}
}

%% memo:
\memoanswer{
Ⅰ．由螺旋测微器与游标卡尺的读数规则可得两者的读数：$d=10+24.3\times0.01=10.243\Umm$，$D=242+4\times0.05=242.20\Umm=24.220\Ucm$。圆盘转动的角速度为 $\omega=\frac{\theta}{t}$，而 $\theta=\frac{d}{2\pi r}\times2\pi$，综合两式并代入数据可得 $\omega=1.69\Urads$。

Ⅱ．（1）测量值为 $R=\frac{U}{I}=800\UO$，因电流表外接，所以 $R=\frac{R_{x}R_{V}}{R_{x}+R_{V}}$，故真实值为 $R_{x}=1000\UO$，对应的百分误差 $\Delta=\frac{1000-800}{1000}=20\%$。

（2）电流表内接时，百分误差 $\Delta'=\left|\frac{1050-1000}{1000}\right|=5\%$。
}

\item
%% number: 23
%% paperName: 2008年四川理综第 23 题 16 分
%% typeId: 6
%% type: 计算题
%% chapter: 直线运动
%% point: 相遇问题
%% method: 无
%% score: 16
%% degree: 450
%% duplicateId: 0
%% body:
A、B 两辆汽车在笔直的公路上同向行驶。当 B 车在 A 车前 $84\Um$ 处时，B 车速度为 $4\Ums$，且正以 $2\Umsq$ 的加速度做匀加速运动；经过一段时间后，B 车加速度突然变为零。A 车一直以 $20\Ums$ 的速度做匀速运动。经过 $12\Us$ 后两车相遇。问 B 车加速行驶的时间是多少？

%% answer:
\jdanswer{$6\Us$}

%% memo:
\memoanswer{
设 A 车的速度为 $v_{A}$，B 车加速行驶时间为 $t$，两车在 $t_{0}$ 时相遇。则有 $s_{A}=v_{A}t_{0}$ ①，$s_{B}=v_{B}t+\frac{1}{2}at^{2}+(v_{B}+at)(t_{0}-t)$ ②。式中 $t_{0}=12\Us$，$s_{A}$、$s_{B}$ 分别为 A、B 两车相遇前行驶的路程。依题意有 $s_{A}=s_{B}+s$ ③，式中 $s=84\Um$。由①②③式得 $t^{2}-2t_{0}t+\frac{2\left[(v_{A}-v_{B})t_{0}-s\right]}{a}=0$ ④。代入题给数据 $v_{A}=20\Ums$，$v_{B}=4\Ums$，$a=2\Umsq$，有 $t^{2}-24t+108=0$ ⑤。式中 $t$ 的单位为 $\Us$。解得 $t_{1}=6\Us$，$t_{2}=18\Us$ ⑥。$t_{2}=18\Us$ 不合题意，舍去。因此，B 车加速行驶的时间为 $6\Us$。
}

\item
%% number: 24
%% paperName: 2008年四川理综第 24 题 19 分
%% typeId: 6
%% type: 计算题
%% chapter: 磁场
%% point: 带电粒子在磁场中的运动
%% method: 无
%% score: 19
%% degree: 450
%% duplicateId: 0
%% body:
如图，一半径为 $R$ 的光滑绝缘半球面开口向下，固定在水平面上。整个空间存在匀强磁场，磁感应强度方向竖直向下。一电荷量为 $q$（$q>0$）、质量为 $m$ 的小球 P 在球面上做水平的匀速圆周运动，圆心为 $O'$。球心 O 到该圆周上任一点的连线与竖直方向的夹角为 $\theta$（$0<\theta<\frac{\pi}{2}$）。为了使小球能够在该圆周上运动，求磁感应强度大小的最小值及小球 P 相应的速率。重力加速度为 $g$。

\onepicture[w=5.0cm, align=right]{figs/24.png}

%% answer:
\jdanswer{$B_{\min}=\frac{2m}{q}\sqrt{\frac{g}{R\cos\theta}}$，$v=\sqrt{\frac{gR}{\cos\theta}}\sin\theta$}

%% memo:
\memoanswer{
据题意，小球 P 在球面上做水平的匀速圆周运动，该圆周的圆心为 $O'$。P 受到向下的重力 $mg$、球面对它沿 OP 方向的支持力 $N$ 和磁场的洛伦兹力 $f=qvB$ ①，式中 $v$ 为小球运动的速率。洛伦兹力 $f$ 的方向指向 $O'$。根据牛顿第二定律 $N\cos\theta-mg=0$ ②，$f-N\sin\theta=m\frac{v^{2}}{R\sin\theta}$ ③。由①②③式得 $v^{2}-\frac{qBR\sin\theta}{m}v+\frac{qR\sin^{2}\theta}{\cos\theta}=0$ ④。由于 $v$ 是实数，必须满足 $\Delta=\left(\frac{qBR\sin\theta}{m}\right)^{2}-\frac{4gR\sin^{2}\theta}{\cos\theta}\geqslant0$ ⑤。由此得 $B\geqslant\frac{2m}{q}\sqrt{\frac{g}{R\cos\theta}}$ ⑥。可见，为了使小球能够在该圆周上运动，磁感应强度大小的最小值为 $B_{\min}=\frac{2m}{q}\sqrt{\frac{g}{R\cos\theta}}$ ⑦。此时，带电小球做匀速圆周运动的速率为 $v=\frac{qB_{\min}R\sin\theta}{2m}$ ⑧。由⑦⑧式得 $v=\sqrt{\frac{gR}{\cos\theta}}\sin\theta$ ⑨。
}

\item
%% number: 25
%% paperName: 2008年四川理综第 25 题 20 分
%% typeId: 6
%% type: 计算题
%% chapter: 运动和力
%% point: 动量和冲量
%% method: 无
%% score: 20
%% degree: 450
%% duplicateId: 0
%% body:
一倾角为 $\theta=45^{\circ}$ 的斜面固定于地面，斜面顶端离地面的高度 $h_{0}=1\Um$，斜面底端有一垂直于斜面的固定挡板。在斜面顶端自由释放一质量 $m=0.09\Ukg$ 的小物块（视为质点）。小物块与斜面之间的动摩擦因数 $\mu=0.2$。当小物块与挡板碰撞后，将以原速返回。重力加速度 $g=10\Umsq$。在小物块与挡板的前 4 次碰撞过程中，挡板给予小物块的总冲量是多少？

\onepicture[w=4.5cm, align=right]{figs/25.png}

%% answer:
\jdanswer{$I=0.43(3+\sqrt{6})\UNs$}

%% memo:
\memoanswer{
解法一：设小物块从高为 $h$ 处由静止开始沿斜面向下运动，到达斜面底端时速度为 $v$。由功能关系得 $mgh=\frac{1}{2}mv^{2}+\mu mg\cos\theta\frac{h}{\sin\theta}$ ①。以沿斜面向上为动量的正方向，按动量定理，碰撞过程中挡板给小物块的冲量 $I=mv-m(-v)$ ②。设碰撞后小物块所能达到的最大高度为 $h'$，则 $\frac{1}{2}mv^{2}=mgh'+\mu mg\cos\theta\frac{h'}{\sin\theta}$ ③。同理，有 $mgh'=\frac{1}{2}mv'^{2}+\mu mg\cos\theta\frac{h'}{\sin\theta}$ ④，$I'=mv'-m(-v')$ ⑤。式中 $v'$ 为小物块再次到达斜面底端时的速度，$I'$ 为再次碰撞过程中挡板给小物块的冲量。由①②③④⑤式得 $I'=kI$ ⑥，式中 $k=\sqrt{\frac{\tan\theta-\mu}{\tan\theta+\mu}}$ ⑦。由此可知，小物块前 4 次与挡板碰撞所获得的冲量成等比级数，首项为 $I_{1}=2m\sqrt{2gh_{0}(1-\mu\cot\theta)}$ ⑧。总冲量为 $I=I_{1}+I_{2}+I_{3}+I_{4}=I_{1}(1+k+k^{2}+k^{3})$ ⑨。由 $1+k+k^{2}+\cdots+k^{n-1}=\frac{1-k^{n}}{1-k}$ ⑩，得 $I=\frac{1-k^{4}}{1-k}\cdot2m\sqrt{2gh_{0}(1-\mu\cot\theta)}$（11）。代入数据得 $I=0.43(3+\sqrt{6})\UNs$（12）。

解法二：设小物块从高为 $h$ 处由静止开始沿斜面向下运动，小物块受到重力、斜面对它的摩擦力和支持力，小物块向下运动的加速度为 $a$，依牛顿第二定律得 $mg\sin\theta-\mu mg\cos\theta=ma$ ①。设小物块与挡板碰撞前的速度为 $v$，则 $v^{2}=2a\frac{h}{\sin\theta}$ ②。以沿斜面向上为动量的正方向，按动量定理，碰撞过程中挡板给小物块的冲量为 $I=mv-m(-v)$ ③。由①②③式得 $I_{1}=2m\sqrt{2gh(1-\mu\cot\theta)}$ ④。设小物块碰撞后沿斜面向上运动的加速度大小为 $a'$，依牛顿第二定律有 $mg\sin\theta+\mu mg\cos\theta=ma'$ ⑤。小物块沿斜面向上运动的最大高度为 $h'=\frac{v^{2}}{2a'}\sin\theta$ ⑥。由②⑤⑥式得 $h'=k^{2}h$ ⑦，式中 $k=\sqrt{\frac{\tan\theta-\mu}{\tan\theta+\mu}}$ ⑧。同理，小物块再次与挡板碰撞所获得的冲量 $I'=2m\sqrt{2gh'}(1-\mu\cot\theta)$ ⑨。由④⑦⑨式得 $I'=kI$ ⑩。由此可知，小物块前 4 次与挡板碰撞所获得的冲量成等比级数，首项为 $I_{1}=2m\sqrt{2gh_{0}(1-\mu\cot\theta)}$（11）。总冲量为 $I=I_{1}+I_{2}+I_{3}+I_{4}=I_{1}(1+k+k^{2}+k^{3})$（12）。由 $1+k+k^{2}+\cdots+k^{n-1}=\frac{1-k^{n}}{1-k}$（13），得 $I=\frac{1-k^{4}}{1-k}\cdot2m\sqrt{2gh_{0}(1-\mu\cot\theta)}$（14）。代入数据得 $I=0.43(3+\sqrt{6})\UNs$（15）。
}

\end{enumerate}

\end{document}
